% TOP_PROBLEM: {"id": 30006753, "title": "Global Regularity of Solutions to the Inviscid Surface Quasi-Geostrophic Equation", "category_id": 9, "category": {"id": 9, "name": "pde", "display_name": "Partial Differential Equations", "description": "PDEs and their applications in physics and geometry.", "slug": "pde", "order_index": 9, "created_at": "2026-07-31T15:26:25.671Z"}, "status": "open"}
% FIRST_POSTED: null
% !TeX program = lualatex
% Compile twice: lualatex 30006753.tex
% All references and prose are embedded; no BibTeX database or external inputs.
\documentclass[11pt,a4paper]{article}
\usepackage[margin=24mm,headheight=14pt]{geometry}
\usepackage{fontspec}
\setmainfont{Latin Modern Roman}
\setsansfont{Latin Modern Sans}
\setmonofont{Latin Modern Mono}
\usepackage{amsmath,amssymb,mathtools}
\usepackage{unicode-math}
\setmathfont{Latin Modern Math}
\usepackage{microtype}
\usepackage{enumitem,xurl,needspace}
\usepackage[dvipsnames]{xcolor}
\usepackage{hyperref}
\hypersetup{unicode=true,colorlinks=true,linkcolor=MidnightBlue,urlcolor=MidnightBlue,
 pdftitle={Research attempt: Problem 30006753},pdfauthor={Research notebook},bookmarksnumbered=true}
\usepackage{bookmark}
\usepackage{tocloft}
\setlength{\cftsecnumwidth}{3em}
\cftsetpnumwidth{2.5em}
\cftsetrmarg{4em}
\usepackage{titlesec}
\titleformat{\section}{\Large\bfseries\raggedright}{\thesection}{0.7em}{}
\usepackage{fancyhdr}
\pagestyle{fancy}\fancyhf{}
\fancyhead[L]{\small\sffamily Open Problems Math | Research notebook}
\fancyhead[R]{\small Problem 30006753 | 27 September 2026}
\fancyfoot[C]{\thepage}
\setlength{\parindent}{0pt}
\setlength{\parskip}{3pt plus 1pt}
\setlength{\emergencystretch}{3em}
\setcounter{tocdepth}{1}
\setcounter{secnumdepth}{1}
\setlist[itemize]{leftmargin=3em,itemsep=5pt,topsep=4pt,parsep=0pt}
\allowdisplaybreaks
\newcommand{\N}{\mathbb{N}}
\newcommand{\Z}{\mathbb{Z}}
\newcommand{\Q}{\mathbb{Q}}
\newcommand{\R}{\mathbb{R}}
\newcommand{\C}{\mathbb{C}}
\newcommand{\F}{\mathbb{F}}
\newcommand{\A}{\mathbb{A}}
\newcommand{\PP}{\mathbb P}
\newcommand{\E}{\mathbb{E}}
\newcommand{\eps}{\varepsilon}
\newcommand{\dd}{\,\mathrm d}
\newcommand{\sref}[2]{\hyperlink{source:#1:#2}{[S#2]}}
\newcommand{\eref}[2]{\hyperlink{extra:#1:#2}{[E#2]}}
% Turn the next switch off for a shorter reading copy. The quotations preserve
% every exactTarget field, including qualifications omitted from a short summary.
\newif\ifshowcatalogue
\showcataloguetrue
\newcommand{\cataloguescope}[1]{\ifshowcatalogue
 \paragraph{Exact scope in the supplied catalogue.}
 \begin{quote}\small #1\end{quote}\fi}
\usepackage{etoolbox}
\AtBeginEnvironment{itemize}{\small}
\numberwithin{equation}{section}
% Standalone export. Original research content preserved.
\begin{document}
\setcounter{section}{0}
% ================================================================
% problemId: 30006753
% Canonical problem ID: 30006753
% exactTarget SHA-256: a50ff32be14b69e5bf6cef4381e8503552eb345abd09fda4b65b9f5ac5d803bd
\clearpage
\section*{\texorpdfstring{Global Regularity of Solutions to the Inviscid Surface Quasi-Geostrophic Equation}{Global Regularity of Solutions to the Inviscid Surface Quasi-Geostrophic Equation}}\label{30006753}
\subsection{Definitions and mathematical statement}
For the scalar $\theta(t,x)$ on $\R^2$, let $R_j$ be the Riesz transform with Fourier multiplier $-i\xi_j/|\xi|$ and set $R^\perp\theta=(-R_2\theta,R_1\theta)$. Inviscid SQG is
\[
 \partial_t\theta+u\cdot\nabla\theta=0,\qquad u=R^\perp\theta,
 \qquad\theta(0,x)=\theta_0(x).
\]
The problem is global smooth existence for every permitted smooth decaying finite-energy datum, or a finite-time breakdown for one such datum. ``Inviscid'' means no dissipative fractional-Laplacian term is present. The precise initial norm and smoothness class should follow the cited formulation rather than a weaker low-regularity ill-posedness result.
\par\smallskip\textit{Formulation and concept references:} \sref{30006753}{1}.
\cataloguescope{For the inviscid SQG equation partial\_t theta + u dot grad theta = 0 on R\textasciicircum{}2, with u = R-perp theta, prove global smooth existence for every sufficiently smooth decaying finite-energy initial datum, or construct such a datum whose solution develops a finite-time singularity.}
\subsection{Short English statement}
Can a smooth inviscid SQG flow form a singularity in finite time, or must every allowed smooth initial state remain smooth forever?
% BEGIN RESEARCH ADDITION 148
\subsection{Research attempt}

\paragraph{Current frontier and data class.}
The supplied paper's smooth norm-inflation examples remain smooth on their
stated interval; its continuous-loss theorem starts in $H^s$, $3/2<s<2$,
not in the smooth class of this target. These results cannot be relabeled
finite-time blow-up from one smooth finite-energy datum.
\sref{30006753}{1}\eref{30006753}{3}
Hou--Qin--Sire--Wu's 2026 construction deliberately permits infinite energy
and a one-dimensional reduction; its ansatz is not a smooth decaying
finite-energy counterexample here. \eref{30006753}{1}
Nontrivial global smooth solution families are known, but do not establish
the every-datum assertion. \eref{30006753}{2}

\paragraph{Attack route.}
A continuation criterion suggests controlling time-integrated gradient
stretching. A frequency-space alternative exploits an exact cancellation
when interacting frequencies have equal length. We derive that cancellation,
quantify its gain in a thin annulus, and actively test whether the annulus
is preserved. The resulting leakage shows why the initial estimate alone
cannot close a global bootstrap. The deductions below are explicit
calculations, not a claim of a new regularity theorem.

\paragraph{Attempt: expose the continuation quantity.}
For a smooth decaying solution and integer $s\ge3$, let
$X=\|\theta\|_{H^s}$, $M=\|\nabla\theta\|_\infty$, and
$\Theta_2=\|\theta_0\|_2$. Transport by the divergence-free velocity
conserves $\Theta_2$. A dyadic decomposition gives, for $J\ge1$,
\[
 \|\nabla u\|_\infty
 \le C\Theta_2+CJM+C_s2^{-(s-2)J}X.
\]
Indeed low frequencies are bounded using the $L^2$ norm; each middle block
is a localized order-zero multiplier applied to $\nabla\theta$, with
uniformly bounded convolution $L^1$ kernel. At high frequencies, Bernstein
in two dimensions gives
$\|\nabla u_j\|_\infty\le C2^{2j}\|\theta_j\|_2$, and Cauchy--Schwarz
sums the geometric tail using $s>2$.
Choosing $J=\lceil\log_2(e+X)/(s-2)\rceil$ yields
\begin{equation}\label{30006753:log}
 \|\nabla u\|_\infty
 \le C_s\bigl(1+\Theta_2+M\log(e+X)\bigr).
\end{equation}

The differentiated transport energy estimate is
\[
 X'\le C_s\bigl(\|\nabla u\|_\infty+M\bigr)X.
\]
Here the leading transport term cancels by $\nabla\cdot u=0$; the
Leibniz commutator terms are bounded by
$C_s(\|\nabla u\|_\infty\|\theta\|_{H^s}
+\|u\|_{H^s}\|\nabla\theta\|_\infty)$, using product interpolation
for intermediate derivative orders. The Riesz multiplier is bounded on
$H^s$, so $\|u\|_{H^s}\le CX$. With $W=\log(e+X)\ge1$,
\eqref{30006753:log} therefore implies
\[
 W'\le C_s(1+\Theta_2+M)W.
\]
Consequently $\int_0^TM(t)\,dt<\infty$ bounds every initially finite
$H^s$ norm and gives continuation by the usual local smooth theory. This
proves the familiar sufficient criterion rather than deriving its unknown
hypothesis from conserved scalar norms.

For $g=\nabla\theta\ne0$ and $\xi=g/|g|$, the exact material identity is
\[
 (\partial_t+u\cdot\nabla)g=-(Du)^{\mathsf T}g,\qquad
 (\partial_t+u\cdot\nabla)\log|g|
       =-\xi^{\mathsf T}(Du)^{\mathsf T}\xi.
\]
There is no pointwise sign in this stretching term. Controlling its
positive accumulation along trajectories is a possible geometric route,
but incompressibility alone supplies no such bound.

\paragraph{A thin-shell cancellation estimate.}
Use $\widehat\theta(\xi)=\int e^{-ix\cdot\xi}\theta(x)\,dx$ and
$\xi^\perp=(-\xi_2,\xi_1)$. The exact velocity convention in the question
gives $\widehat u(\xi)=-i\xi^\perp|\xi|^{-1}\widehat\theta(\xi)$.
Symmetrizing the Fourier convolution yields
\begin{equation}\label{30006753:kernel}
 \widehat{u\cdot\nabla\theta}(q)
 =\frac{1}{2(2\pi)^2}\int_{k+\ell=q}
 (k^\perp\cdot\ell)
 \left(\frac1{|k|}-\frac1{|\ell|}\right)
 \widehat\theta(k)\widehat\theta(\ell)\,dk.
\end{equation}
The factor vanishes exactly when the two radii agree, as well as for
parallel interacting frequencies. Suppose at a single time
\[
 \operatorname{supp}\widehat\theta\subset
 \{R\le|\xi|\le(1+\epsilon)R\},
 \qquad R\ge1,\quad0<\epsilon\le\tfrac12.
\]
The symmetrized multiplier has magnitude at most $C\epsilon R$, and its
output support has radius at most $3R$. Young's convolution inequality and
Cauchy--Schwarz on the annulus then prove
\begin{align}
 \|u\cdot\nabla\theta\|_{H^s}
 &\le C_s\epsilon R(1+R)^s
                   \|\widehat\theta\|_1\|\theta\|_2\notag\\
 &\le C_s\epsilon^{3/2}R^2(1+R)^s\|\theta\|_2^2.
 \label{30006753:shell}
\end{align}
The second step uses annulus area $O(\epsilon R^2)$. This is a proved
instantaneous estimate, not an assumption of random phases or a
conservation law for the spectral support.

\paragraph{Stress test: the annulus leaks immediately.}
Take $0<\epsilon<1/4$ and frequencies
\[
 k=R(1+\epsilon/4,0),\qquad
 \ell=R(0,1+3\epsilon/4).
\]
They lie strictly within the annulus, but $|k+\ell|>(1+\epsilon)R$.
The symmetrized multiplier in \eqref{30006753:kernel}, including its factor
$1/2$ but excluding $(2\pi)^{-2}$, is $\epsilon R/4\ne0$.
Thus orthogonal frequencies of slightly different radii generate a frequency
outside the initial annulus.

This can be realized with smooth finite-energy data, not just plane waves.
Choose nonnegative smooth Fourier bumps with small disjoint supports around
$\pm k,\pm\ell$, symmetric under $\xi\mapsto-\xi$. Their inverse transform
is real and Schwartz. At frequencies near $k+\ell$, only the two cross
interactions contribute; their symmetrized multiplier has the same strict
sign on sufficiently small supports. The convolution is consequently
nonzero on an open output set outside the annulus. The local smooth solution
is differentiable in a Sobolev space, so projection to that output set gives
$\widehat\theta(t)=t\,\partial_t\widehat\theta(0)+o(t)$ there in $L^2$.
Its spectral support has left the annulus for sufficiently small positive
time. Equation~\eqref{30006753:shell} cannot therefore be iterated by assuming
the same support restriction.

Exact support on one circle is not a nontrivial finite-energy escape from
this difficulty: an $L^2(\R^2)$ Fourier transform supported on a
measure-zero circle is zero. On a torus monochromatic modes may be
nontrivial and stationary, but changing domain would change the target.
Radial Schwartz data on $\R^2$ do give stationary solutions, since
$u=-\nabla^\perp\Lambda^{-1}\theta$ is tangent to circles and
$\nabla\theta$ is radial. That special case does not control angular
perturbations or arbitrary data.

Finally $\theta_\lambda(t,x)=\theta(\lambda t,\lambda x)$ is the
amplitude-preserving scaling. Initially it keeps $\|\theta\|_\infty$
fixed, decreases $\|\theta\|_2$ by $\lambda^{-1}$, and increases
$\|\nabla\theta\|_\infty$ by $\lambda$. Bounded conserved upper norms
therefore cannot give an a priori gradient bound without additional
structure. This is a scaling countertest, not a singular solution.

\paragraph{Outcome.}
\textbf{Outcome: Exploratory approach with unresolved gap.}
The attempt supplies the continuation calculation, an exact equal-radius
cancellation and a quantitative thin-annulus bound. It also proves that
smooth finite-energy data can immediately generate frequencies outside that
annulus, invalidating the proposed invariant-shell bootstrap. The missing
step is control of the newly generated spectrum or of integrated gradient
stretching for arbitrary smooth data. Neither a global bound nor a
finite-time singularity in the target class is established.

% END RESEARCH ADDITION 148
\subsection{Sources}
\begin{itemize}\raggedright
\item[\textnormal{[S1]}]\hypertarget{source:30006753:1}{} Cordoba, Martinez-Zoroa, and Ozanski, Instantaneous Continuous Loss of Regularity for the SQG Equation, Introduction.
\url{https://www.sciencedirect.com/science/article/pii/S0001870825004517}.
{\small\textit{Catalogue formulation source.}}
% BEGIN RESEARCH REFERENCES 148
\item[\textnormal{[E1]}]\hypertarget{extra:30006753:1}{} T. Y. Hou, X. Qin,
Y. Sire and Y. Wu, \emph{Self-similar blow-up profile for the one-dimensional
reduction of generalized SQG with infinite energy}, arXiv:2603.11301v1,
11 March 2026, Introduction and Theorem 1.1.
\url{https://arxiv.org/html/2603.11301v1}.
\item[\textnormal{[E2]}]\hypertarget{extra:30006753:2}{} A. Castro, D. C\'ordoba
and J. G\'omez-Serrano, \emph{Global smooth solutions for the inviscid SQG
equation}, Memoirs of the AMS 266 (2020), no. 1292; arXiv:1603.03325,
abstract and solution-family theorem. \url{https://arxiv.org/abs/1603.03325}.
\item[\textnormal{[E3]}]\hypertarget{extra:30006753:3}{} D. C\'ordoba,
L. Mart\'inez-Zoroa and W. S. O\.{z}a\'nski, \emph{Instantaneous continuous
loss of regularity for the SQG equation}, accessible manuscript
arXiv:2409.18900v1, Section 1.1, Theorem 1 and continuous-loss statement;
manuscript of [S1]. \url{https://arxiv.org/html/2409.18900v1}.

% END RESEARCH REFERENCES 148
\end{itemize}

\end{document}
